% chapters/09_titans_and_multi_timescale_memory.tex
% 第9章：Titans 与多时间尺度记忆
\chapter{Titans 与多时间尺度记忆：别把鸡蛋放在一个篮子里}
\label{ch:titans}

\epigraph{``You cannot expect a single neural network layer to handle both the immediately relevant context and the lifelong accumulation of facts. Memory must be differentiated by its timescale.''}{--- Titans: Learning to Memorize}

\section{本章想回答什么问题}

在第 \ref{ch:ttt_layers} 章中，我们将“大显存上下文硬搜（KV Cache）”优化成了“带有推断层微调的小型模型（TTT-MLP）”。它利用自己的内部层网络容纳并在新信息经过时做适应调节更新，容量与非线性泛化性极强。

但这又带来了一个极其严重且客观存在的死局：\textbf{灾难性遗忘（Catastrophic Forgetting）}。
\begin{quote}
\textit{既然这个微型网络是一个整体，那么如果在100万个字符的处理长河中，每时每刻的新字都在逼迫它沿损失方向微调权重去记忆当前的事物。到了文章末尾，它最早期记住的信息哪怕是最隐秘非线性的深层权重，不也早被后来无数不相干的高频次普通更新覆盖重置而面目全非了吗？}
\end{quote}

本章将聚焦于在超长文本（百万量级以上 token）时代的最新演进：Titans（Learning to Memorize at Test Time）。我们将回答如果面对极限超长的信息流任务需要兼顾长短记忆保留，我们应该如何在架构中分立“短时精确焦点”与“长篇容纳存储器”，并用一个叫做惊讶度（Surprise）的关键机制来做选择性的高质量遗忘与吸收。

\section{最小例子：沙海淘金的极端过滤考验}
\label{sec:titans_minimal}

\begin{toybox}[title=玩具问题 F：千万级 Token 的“草垛寻针”]
试想一种任务：在一部上千万字的长篇废话巨著（草垛）的第 100 页里随机塞入了一句：“特别任务暗号：红狐狸跳墙。”除此之外一千万字全是干扰性的废料。最后一句直接向模型提问：“暗号是什么？”（著名的 Needle In A Haystack 草垛寻针基准）。

\textbf{如果在纯标准 TTT 架构下：}
刚读完暗号时，TTT隐性层进行了更新把它内化到了权重里。但是接着的一千万字废话犹如不休不止的海浪，每一次都产生更新的冲刷，不断强行微调这个薄弱隐层矩阵使其努力去适应和拟合最新的那点废料。最后由于权重飘移漂白，模型彻底把暗号“覆盖抹除”了。

\textbf{Titans 架构下：}
当废话到来时，Titans中独有的由惊诧度（Surprise）算子构成的一套监察守门系统发现“废话”极度符合普遍预期，没有信息量（Surprise $\to 0$），于是不允许把它们进行大的写入（控制记忆网络参数封印锁定不改）！而只把它们暂时维持在并行短存的小框里处理完迅速丢弃。而当“暗号”出现这种反常且不匹配平时预期的强信息突刺时，惊讶度飙升，大开闸门重重地将其向负责长期记忆的网络中烙印！这就极大免疫了由于无差别连续调整带来的干扰摧毁。从而即使在百万字后，也能把深藏不漏的针给调出来。
\end{toybox}

\section{直觉解释：将草稿纸与刻表分类}
\label{sec:titans_intuition}

在模型中强加一个庞然大物般的无差别学习器是不对的。正确的是进行明确的分工协作计算系统。

Titans 框架抛弃了“用一个模块打天下”的想法。它并立（或层配）了两种截然不同的运算资源池组和提取管线：
\begin{enumerate}
    \item \textbf{并行窗口注意力模块（也就是草稿纸）：}只处理当前最新的一小段，极为精确且完全不遗忘，算完就换页丢弃。这用于承担对即时语境的高度词法结构依赖操作。
    \item \textbf{神经常期记忆模块（那是雕刻泥板）：}这是一个带有持久累积效应的参数学习层。为了不对这块珍贵容量做胡乱消磨，加装了基于“当前记忆的偏差预期程度”的判断标量来作为进入这个大网更新其可塑性的倍差阀门！
\end{enumerate}

最终输出是被这两个途径分别算完并根据情况以不同权比（Alpha门控）进行捏合汇出给下游输出决定的。

\begin{neurosciencebox}[title=神经科学直觉：系统融合与长短时海马体区分]
\textbf{计算功能类比：}
Titans 这种区分短视窗计算和极其慎重的永久写录更新极其巧妙地类比了神经学中\textbf{短时工作记忆向长期皮层固化存储}（Cortical Consolidation）的转移分工。
\begin{itemize}
    \item 局部精确的“窗口注意力（Window Attention）”可以类比为前额叶主导的**工作记忆（Working Memory）**。在这个区域里，你非常清楚对方上一分钟前说了什么的精确词句排序；但你过两天就全忘了，它容量固定但极其精准且抛弃极快。
    \item 而带着筛选器通过惊奇重塑来进行网络权重刻录的模块，非常贴切地展示了\textbf{海马结构与杏仁核联合（Hippocampus \& Amygdala interplay）对强情绪/强异常事件发生极其深度的长期突触修改强留（LTP）}！只有引起极大意外（Surprise）的事物或者不断由于高度显著而击穿注意阀的值，神经微系统才允许为它消耗那极度高昂代谢并永久拉动全局突触链接（改变皮层网络的重量分布用以伴随终身提取）。因为你的大脑也是一再规避灾难性突触同质重塑失效的！
\end{itemize}
\end{neurosciencebox}

\begin{analogyboundarybox}[title=这并不是说“模型就是大脑”]
虽然从长/短分流以及“异变情绪（惊叹门）”把关记忆写入的做法听起来就像是对认知心理学的完全模仿，但这其中有着绝不容混淆的形式边界：
\begin{itemize}
    \item \textbf{更新规则的严酷冰冷：}在模型中，Surprise 被物理地定性为了 $s_t = \| M_{\Phi} - v \|$ 一个刚性的误差数学标量乘积，没有任何类似于神经系统的生化传递阻滞和多路信号非线性突现特征。它只是一场算计好的数学梯度加权步长惩罚控制。
    \item \textbf{结构并非绝对实体隔离：}Titans 的两个模块实际上同处一个大型可导架构下并在一个前向中进行紧耦合向量拼接加法运算；而在脑科学中，海马长期化是在睡眠与离线过程大量伴有电活动回放（Replay）并跨越大脑叶层发生的物理电传导改建，它涉及跨组织的宏大空间漫游生理变性支持而不是并排简单的输出相加权。永远把这种机制视作在数学系统上\textbf{极具高阶抽象概括和阉割重组过后的优化效能手段而已}！
\end{itemize}
\end{analogyboundarybox}

\section{数学形式化：惊诧度（Surprise）作为可导滤筛}
\label{sec:titans_math}

我们展示 Titans 之所以精妙的关键目标公式，即长期的“神经记忆模块”核心计算如何兼顾自我保护与自我修正。

\subsection{惊诧门控与状态修正步进}
假设 $\Phi_{t-1}$ 是这个大记忆库此刻带有的网络参数。当外部慢权重投射入新的查询-输出对比对（$k_t, v_t$）。
网络此时会通过读取其自身对 $k_t$ 的既有理解，产生输出：
\begin{equation}
   \hat{v}_t = M(k_t; \Phi_{t-1})
\end{equation}

在经典的隐式梯度步（如上一章的 TTT 更新那般），参数调整会顺着 $\eta \nabla_{\Phi} \ell$ 继续发生大偏移。

而这儿我们引入“惊诧度”，在数学里，它被自然定义为当期输出和预期值的真实向量误差差值的直接标量范数衡量项：
\begin{equation}
    s_t = \sigma( \| \hat{v}_t - v_t \|_2 - b ) \quad （或者配合其他线性缩放）
\end{equation}
此处 $\sigma$ 表示通过约束后的系数（或甚至就是包含在梯度中的那个缩放比），$b$ 做基线抵消。这使得当预测已经挺准确（代表已经完全被同化和记住这类知识）时，或者极其微小偏移时，算式产生极其逼近零的正向增益。

将其与梯度下降混合：
\begin{equation}
    \Phi_t = \Phi_{t-1} - \eta \cdot s_t \cdot \nabla_{\Phi} \ell( \hat{v}_t , v_t )  \ + \ \alpha\Delta\Phi_{\text{momentum}}
\end{equation}

通过这一个控制项 $s_t$，在梯度图反拉时如果误差微弱，则强行以零项挂起，\textbf{强制锁死了对于当前这种无效常见平滑信息的改写响应度门阀！}从而为存储记忆空间进行了巨大减负保全。

\subsection{输出合成组网}
最终预测结果将被结合上述模块加上另一个平行的短距离纯自注意力缓存头（$H_{SA}$）所产生的结果拼装并重投影进行下游任务应用。这种分离即在保留近期 $W$ 视窗的绝对无损下，赋予了超越数个世纪级别时空跳跃的隐式过滤检索能。

\section{代表论文精读与发展}
\label{sec:titans_papers}

\paragraph{\citet{behrouz2024titans} — Titans: Learning to Memorize at Test Time}
这篇重磅之作宣告在抛除极其占有显存的情况下使得真正具有无尽大记忆能力的测试自适应生成成为现实基理底座方案突破。它的主论文并不是空对空讲结构，而是非常实在的放开到近数百万长的寻针等评测级进行消融破坏。证实了仅使用其中一个模块会立刻崩解并在融合状态下展现出了即使对于 $2M$ Token长度也不发生梯度塌陷。

\section{和前后章节的关系}
\label{sec:titans_bridges}

\begin{itemize}
    \item \textbf{对第 \ref{ch:ttt_layers} 章 TTT的打补丁修正：}它是直接作为针对 TTT 如果无边界地吞噬和被改写会导致全盘“参数遗忘毁天灭地”所做的最后一道防御屏障和救生衣架构方案加入。
    \item \textbf{通向第 \ref{ch:unified} 章 和第 \ref{ch:nested} 章的前置线索：}这种基于“当前偏差判断”再去执行“过滤”或“不同更新深度”的手法。本质上说明，在在线持续回归系统里，所有的新数据权重并不是等同处理的；也启示了如果把这套把控时间截断与门控的逻辑泛推，整个深度的结构实质也就是套娃，指引全书理论最高大收尾。
\end{itemize}

\begin{labbox}[title=Lab 9: 设计惊奇阈值的抗干扰遗忘截断效果观察]
\textbf{目标：}利用一个微型的循环回归参数网络，手动加入/去除 Surprise 项，比较他们在长期数据被覆盖中“寻回最初高频重点”的能力。

\begin{enumerate}
    \item \textbf{构造数据流}：总长 $T=10000$ 的 1 维流序列，在前 10 步放入目标 $(0.5, 0.99)$ 这个对。在那后面的 $9990$ 步全部是用一个微小高斯分布震荡于均值附近波动的底噪噪音对。
    \item \textbf{朴素微更新（无把关）}：开启基于 MSE 梯度的普通递推，在走完 $10000$ 步之后去问如果在 $(0.5)$ 此位置它能给出什么解？（可以预见：由于全域不断被细小的高斯噪音对在微调着，最终它会失去对那个大极值 0.99 的精确匹配而是掉到了平庸均值附近）。
    \item \textbf{建立门控（加惊诧判定控制阀）}：设定 $\Delta W = \text{Error} \cdot (\text{如果Error}>0.1 \rightarrow 倍率开大 ；反之倍率切分为 1/1000 衰减)。$
    \item 重新再过一次！再去拿 0.5 测！由于惊诧门直接封挡了后方那些只有 0.01 震荡的连续干扰噪音废数据的污染，参数极高地保持着纯净状态并在大步重写入期间维持住了坚挺响应度！
\end{enumerate}
这就是对系统性知识“被动抵抗性记忆保护”的数学粗体验。
\end{labbox}

\begin{misconceptionbox}
\textbf{常见误解 1：加了门限就能无限大跨度保持不忘所以这个记忆存储模块可以无限用下去不用管。}

\textbf{纠正：}这是一种错觉与过激乐观。这依然受到由于物理矩阵特征降秩极点的压制。如果出现的“极其令人惊奇和具有大信息的知识点”足够多并充斥在整个时间段而不是被那些容易隔离废料噪音包围，那么惊诧滤门会一直处于“强开启态”，最后仍就会出现系统由于权重冲刷带来的交叉退化灾难性溃散和交织记忆幻觉，这就是为什么工业化需要定期的进行跨阶段冷存储洗牌或外部检索重构辅助（如 RAG 系统）。
\end{misconceptionbox}

\begin{misconceptionbox}
\textbf{常见误解 2：有了长时记忆这块，外层的结构就不再起到“检索”的作用而只要全部去存就行了。}

\textbf{纠正：}外全局固定的大架构在此依然极为致命重要。惊奇度的高低从何决定？完全依赖于大架构事先预训练出来的“预判表征函数映射器”。预训练学到了什么是世界本该合理存在的普通运行轨迹，它才会觉得那句废话低分平庸；它才能让门控制阀工作。如果整个主骨干本身未被广袤数据雕琢过，它是看不出“反常与无聊”之区别的盲者，那么后面再精巧的一切保护层依然如失效空物。
\end{misconceptionbox}

\section{练习题}
\label{sec:titans_exercises}

\begin{exercise}
\textbf{不同权衡（Trade-off）分配的讨论：}假设我们在工业推行部署一项实时极其关注于毫秒级“动作指令和位置感联”（极度短期关注依赖）兼少量对过去很久的系统启动设置项的回调的机器人动作系统。在 Titans 构架的两个分别被赋权重为 $\alpha_1$（对短工作窗）与 $\alpha_2$（对神经大记忆区）的拼接里。相比于让其作为一个通用的读写长文章 AI，这个针对极高速实时偏重的系统中你会在超参数训练上偏重什么配方趋势为什么？
\end{exercise}

\begin{exercise}
\textbf{长短期机制与梯度的联结矛盾：}在长短时合并系统被提出之前，对于一个只基于隐性更新的大微模型态推导模型，由于其隐性学习带来的步长微推可能与全时反向的后串形成高度纠缠致使多模化训练难以实现融合计算图优化。在目前长短机制采用分别输出然后汇加的粗分离方式后，是否这种由于短线记忆区仅利用纯软线注意因而“可以断开时序链条单独极速矩阵算尽”带来了关于硬件并发支持率上的本质计算瓶颈优势消除的解释？请用硬件分块概念做推演作答。
\end{exercise}

\section{本章小结}
\label{sec:titans_summary}

\begin{itemize}
    \item \textbf{灾难遗忘的解决之药：}本章解释了当模型使用隐态变网络持续改写的方案中随着数据量暴增必然面临的覆写崩塌。指出并提供了基于短视分离抛弃及长效受控强改录两种截然机制分离架构这一破围思路。
    \item \textbf{通过惊诧（误差）设置防线：}不仅要学习还得学会“如何拒绝平庸学习而抵抗消磨”。Surprise门阀作为一个数学和认知心理上的同频共振结合项展示出了巨大价值，完成了对模型进行记忆纯度提纯保修的关键。
    \item \textbf{融合不代表万灵：}它并不是解决了本质表达的局限，而是极高效地把干扰给截除了。对真正有意义长线交叉检索依然挑战着理论天花板。
    \item \textbf{在整个学习阶梯的登顶前奏：}从解决微调、变身为网络再到增加记忆把控，这条以适应度为追求的设计线已经近乎推到极限完形态；下一步就是将整个这片版图拔高，统统解释为时空上的回归切片（即最终两章）。
\end{itemize}

\paragraph{读完本章，你应该能够：}
\begin{itemize}
    \item 为何“用隐状态在极长跨度做记忆全盘维护”是个灾难以及采用多重子策略是自然应对的最优解给出解释原因与机制构成解剖。
    \item 掌握如何在一个普通的一步参数微进公式上加上一个与梯度无关但是取决于输出误差模长的标量因子来进行受控阀门的门控改写控制操作实现方法理念。
    \item 在不将模型真理解为动物生物脑的前提条件下懂得如何利用生物本初对于危险或意外强保留刺激的直觉来更好地搭建或向外解说你的控制逻辑过滤。
\end{itemize}
